% Part: first-order-logic
% Chapter: beyond
% Section: modal-logics

\documentclass[../../../include/open-logic-section]{subfiles}

\begin{document}

\olfileid{fol}{byd}{mod}

\olsection{المنطقات الموجهة}

لننظر في هذه الجملة الشرطية:
\begin{quote}
  إذا كان جيريمي وحده في تلك الغرفة، فهو ثمل وعار ويرقص على الكراسي.
\end{quote}
قد تصدق هذه الشرطية صدقًا ماديًا، ويكون قولها مع ذلك مضللًا.
ذلك أن السامع يفهم منها صلة بين المقدمة والتالية تزيد على مجرد
كذب المقدمة أو صدق التالية. فكأن القائل لا يخبر بصدقها في هذا
العالم وحده، حيث قد تصدق تحصيلًا لأن جيريمي ليس وحده في الغرفة،
بل يزعم أن التالية \emph{كانت ستصدق لو} صدقت المقدمة. وحاصل هذا
الفهم أن الدعوى صادقة في كل عالم «ممكن»، لا في عالمنا فقط؛ أي
إنها صادقة \emph{بالضرورة}، لا صادقة اتفاقًا في هذا العالم بعينه.

ولضبط هذا المعنى من الضرورة وُضع المنطق الموجه. ففي صورته القضوية
نزيد المنطق القضوي المعتاد مؤثر الصندوق: فإذا كانت $!A$~!!a{formula}،
كانت $\Box !A$ كذلك. ومعنى $\Box !A$ على وجه التقريب أن
$!A$ صادقة \emph{بالضرورة}، أي في كل عالم ممكن. وأما
$\Diamond !A$ فيُجعل في الغالب اختصارًا لـ~$\lnot \Box \lnot !A$،
ومعناه أن $!A$ صادقة \emph{بالإمكان}. ولا تختص هذه الزيادة
بالمنطق القضوي؛ فلنا أن نزيد الجهة في منطق المحمولات أيضًا.

وتصلح !!{structure}s كريبكه دلالةً لهذا المنطق؛ بل إن كريبكه
وضعها أولًا من أجله. ولا يلزم هنا الاقتصار على الترتيبات الجزئية:
بل نأخذ مجموعة من «العوالم الممكنة»~${\OLId{latin-looped}{P}}$، وعلاقة ثنائية
$\Atom{R}{x,y}$ تسمى «إمكان الوصول». وتقريب معنى $\Atom{R}{p,q}$
أن العالم~${\OLId{latin-ordinary}{q}}$ لا ينافي العالم~${\OLId{latin-ordinary}{p}}$؛ فإذا كنا «في» العالم~${\OLId{latin-ordinary}{p}}$،
وجب أن نعدّ صورة العالم~${\OLId{latin-ordinary}{q}}$ ممكنة.

وقد يسمى المنطق الموجه «قصديًا»، في مقابلة المنطق «الامتدادي».
فدلالة المنطق الامتدادي، ومنه الكلاسيكي، لا تقصد إلا عالمًا واحدًا،
هو العالم «الفعلي»؛ وأما الدلالة القصدية فتفترض مجالًا من
الموجودات أغنى من ذلك. وليست الضرورة وحدها مما تعبّر عنه الجهة:
بل تتغير قراءة $\Box$ و$\Diamond$ بتغير المقصود، فتمثل بها
تراكيب لغوية أخرى. ومن أمثلة ذلك:
\begin{enumerate}
\item في منطق القابلية للإثبات، يُقرأ $\Box !A$ على أنه «$!A$ قابلة
  للإثبات»، و$\Diamond !A$ على أنه «$!A$ متسقة».
\item في المنطق المعرفي يمكن قراءة $\Box !A$ على أنه «أعرف أن~$!A$»
  أو «أعتقد أن~$!A$».
\item في المنطق الزمني يمكن قراءة $\Box !A$ على أنه «$!A$ صادقة
  دائمًا»، و$\Diamond !A$ على أنه «$!A$ صادقة أحيانًا».
\end{enumerate}

ثم لا بد من بديهيات وقواعد تضبط هذه الجهة. فالنسق~$\Log{S4}$
يجمع بديهيات المنطق القضوي وقواعده المعتادة والبديهيات الآتية:
\begin{align*}
& \Box (!A \lif !B) \lif (\Box !A \lif \Box !B)\\
& \Box !A \lif !A \\
& \Box !A \lif \Box \Box !A
\intertext{ومع القاعدة «استنتج $\Box !A$ من~$!A$». ويضيف
  $\Log{S5}$ البديهية الآتية:}
& \Diamond !A \lif \Box \Diamond !A
\end{align*}
وتختلف البديهيات الملائمة باختلاف التطبيق؛ فـ S5، مثلًا، يعد
عادةً وصفًا لمفهوم الضرورة المنطقية. ومن محاسن هذه الأنساق أنه
يمكن غالبًا تقييد علاقة إمكان الوصول بقيود طبيعية، فتكون
دلالة نسق ال!!{derivation} سليمة وتامة على !!{structure}s كريبكه
المقيدة بها. فالنسق~$\Log{S4}$ يقابل فئة !!{structure}s كريبكه
ذات علاقة الوصول الانعكاسية المتعدية. وأما $\Log{S5}$ فيقابل فئة
!!{structure}s كريبكه ذات علاقة الوصول {\em الكلية}، حيث يُوصل
من كل عالم إلى كل عالم؛ وعندئذ تصدق $\Box !A$ إذا وفقط إذا
صدقت $!A$ في جميع العوالم.

\end{document}
